% ============================================================
%  Paper VIII  --  Working draft
%  Scale-Separated Dyadic Cancellation for PSWF Galerkin Operators
%  Status: working draft (revised v3: gap-axiom + two-scale)
% ============================================================
\documentclass[12pt]{amsart}

% --- packages ---
\usepackage{amsmath,amssymb,amsthm}
\usepackage{mathtools}
\usepackage{hyperref}
\usepackage{enumerate}

% --- theorem environments ---
\newtheorem{theorem}{Theorem}[section]
\newtheorem{proposition}[theorem]{Proposition}
\newtheorem{conjecture}[theorem]{Conjecture}
\newtheorem{corollary}[theorem]{Corollary}
\newtheorem{lemma}[theorem]{Lemma}
\theoremstyle{definition}
\newtheorem{assumption}[theorem]{Assumption}
\newtheorem{definition}[theorem]{Definition}
\theoremstyle{remark}
\newtheorem{remark}[theorem]{Remark}

% --- macros ---
\newcommand{\norm}[1]{\lVert#1\rVert}
\newcommand{\abs}[1]{\lvert#1\rvert}
\newcommand{\lam}[1]{\lambda_{#1}^{(c)}}
\newcommand{\FAi}{F_{\mathrm{Ai}}}
\newcommand{\NSh}{N_{\mathrm{Sh}}}
\newcommand{\Bm}{\mathcal{B}_m}

% ============================================================
\begin{document}
% ============================================================

\title[Scale-Separated Dyadic Cancellation]
      {Scale-Separated Dyadic Cancellation\\
       for PSWF Galerkin Operators:\\
       Two-Scale Structure of the Spectral Transition}

\author{[Author]}
\date{\today}
\maketitle

% ------------------------------------------------------------
\begin{abstract}
We study the eigenvalue transition of the PSWF concentration
operator in the high-bandwidth regime and establish two
classes of results.

\emph{Rigorous.}
Under Proposition~\ref{prop:gap} (uniform spectral gap
lower bound) and Assumption~B-strong \cite{paperVI},
we verify hypotheses~(H1)--(H3) of Paper~VII \cite{paperVII}
via the \emph{Dyadic Separation Principle}:
for $j \in \mathcal{B}_m(i)$, the ratio
$|\lambda_j - \lambda_{j+1}|/|\lambda_i - \lambda_j|$
is bounded by $C \cdot 2^{-m}$ independently of~$c$,
because the factor $(c/2)^{-1/3}$ cancels between
numerator and denominator.
This cancellation is the operative content of the paper
and depends only on monotonicity and a local gap lower bound.

\emph{Empirical.}
We present numerical evidence for a universal scaling law
governing the global transition profile
(Conjecture~\ref{conj:ULW}):
\[
  \lambda_l \approx \tfrac{1}{2}\operatorname{erfc}
  \!\left(\frac{\ln 2\cdot(l-\NSh-\beta(c))}{S(c)}\right),
  \quad S(c) = (\log c/\pi)^{2/3},
\]
with cross-bandwidth residuals of order $2\cdot 10^{-2}$
across $c \in \{50,100,200\}$, decreasing with~$c$.
The centering parameter $\beta(c)$ drifts slowly
($\beta(50)\approx -0.42$, $\beta(200)\approx -0.36$);
its limiting behavior is an open problem.
A full asymptotic derivation of this law remains open.
\end{abstract}

% ============================================================
\section{Introduction and Context}
\label{sec:intro}
% ============================================================

Paper~VII \cite{paperVII} established an abstract dyadic
cancellation theorem (Theorem~2.1) and reduced the
contraction problem to three hypotheses (H1)--(H3).
H2 was stated as Conjecture~6.2 and is reduced here
to Proposition~\ref{prop:gap}.
This paper is a \emph{conditional reduction theorem}:
under Proposition~\ref{prop:gap} and B-strong,
hypotheses (H1)--(H3) are verified and the contraction
estimate follows.

\medskip\noindent
\textbf{Two-scale structure.}
Our analysis reveals a two-layer structure of the
spectral transition:
\begin{itemize}
\item \emph{Combinatorial scale} (Dyadic Separation):
  requires only a uniform gap lower bound at scale
  $(c/2)^{-1/3}$; this is Proposition~\ref{prop:gap},
  Lemma~\ref{lem:F}, and Corollary~\ref{cor:C}.
  This layer is the rigorous core and is independent
  of the specific form of the eigenvalue profile.
\item \emph{Local scale} (Airy / WKB): provides a
  refined description of increments
  $\lambda_l - \lambda_{l+1}$ via Airy functions;
  this is Proposition~\ref{prop:U} and Lemma~\ref{lem:A}.
  It supplies one verification of Proposition~\ref{prop:gap}
  but is not needed for Lemma~\ref{lem:F} or
  Corollary~\ref{cor:C} directly.
\item \emph{Global scale} (Landau--Widom): the eigenvalue
  profile as a whole follows an error-function law with
  logarithmic width $S(c) = (\log c/\pi)^{2/3}$;
  this is Conjecture~\ref{conj:ULW}.
\end{itemize}

\medskip\noindent
\textbf{The Dyadic Separation Principle.}
For $j \in \mathcal{B}_m(i)$:
both $|\lambda_j-\lambda_{j+1}|$ and $|\lambda_i-\lambda_j|$
scale as $c^{-1/3}$, but $|\lambda_i-\lambda_j|$ carries
an extra $2^m$ from $|i-j| \ge 2^m$, so the ratio is
$O(2^{-m})$ independently of $c$.
This mechanism depends only on monotonicity and the
gap lower bound of Proposition~\ref{prop:gap},
not on the fine Airy or erfc structure.

\medskip\noindent
\textbf{What is not proved here.}
Propositions~\ref{prop:gap} and~\ref{prop:U} and B-strong
remain open; see Remarks~\ref{rem:gap_route},
\ref{rem:propU_route}, and~\ref{rem:open}.
Conjecture~\ref{conj:ULW} is supported by strong numerical
evidence but has no analytic proof at present.

% ============================================================
\section{Setup and Notation}
\label{sec:setup}
% ============================================================

Throughout, $c \ge 1$; $\lambda_k \in (0,1)$ ordered
decreasing. $\NSh := 2c/\pi$.

\medskip\noindent
\textbf{Airy zeros and $A_0$.}
All zeros of $\mathrm{Ai}$ are negative;
$a_1 \approx -2.3381$ is the largest \cite[Sec.~9.9]{DLMF}.
Fix $A_0 \in (0,|a_1|)$; then $[-A_0,A_0]$ is zero-free.

\medskip\noindent
\textbf{Compactness constants (depend only on $A_0$):}
\begin{align*}
  c_* &:= \min_{[-A_0,0]}\FAi > 0,\\
  c_{\min} &:= \min_{[-A_0,A_0]}\mathrm{Ai}^2 > 0,\quad
  C_{\max} := \max_{[-A_0,A_0]}\mathrm{Ai}^2 < \infty.
\end{align*}

\medskip\noindent
\textbf{Lipschitz constant:}
$L_{A_0} := \sup_{[-A_0,A_0]}|\frac{d}{dx}\mathrm{Ai}^2|
= 2\sup|\mathrm{Ai}\,\mathrm{Ai}'| < \infty.$

\medskip\noindent
\textbf{Monotonicity fact.}
$\FAi$ is strictly decreasing
($\FAi' = -\mathrm{Ai}^2 \le -c_{\min} < 0$ on $[-A_0,A_0]$).

\medskip\noindent
\textbf{All constants} depend only on $A_0$; none
depend on $c$, $i$, $j$, $m$.

\medskip\noindent
\textbf{Rescaled index:}
$x_l := (l-\NSh)(c/2)^{-1/3}$;
spacing $x_{l+1}-x_l = (c/2)^{-1/3}$ exactly.

\medskip\noindent
\textbf{LW scaling variable:}
$S(c) := (\log c/\pi)^{2/3}$,\quad
$Z_l := \ln 2 \cdot (l - \NSh - \beta(c))/S(c)$.

\medskip\noindent
\textbf{Notation:}
$A_{ij} := P_{ij}c^{1/2}$,\;
$P_{ij} := c|\lambda_i-\lambda_j|/
((1-\lambda_i)(1-\lambda_j))$;\;
$\Bm(i) := \{j: 2^m \le |i-j| < 2^{m+1}\}$.

\medskip\noindent
\textbf{Bound on $1-\lambda_j$.}
Under~(U) with $C_1 c_0^{-2/3} \le c_*/2$:
\begin{equation}\label{eq:onemlam}
  \tfrac{1}{2}\FAi(x_j) \le 1-\lambda_j \le 2\FAi(x_j)
\end{equation}
uniformly in all transition-window indices.

% ============================================================
\section{Three-Layer Spectral Structure}
\label{sec:structure}
% ============================================================

We formulate three propositions at increasing levels of
specificity.
Proposition~\ref{prop:gap} is the axiomatic minimum
required by the rigorous core.
Proposition~\ref{prop:U} is a stronger local statement
supplying one route to Proposition~\ref{prop:gap}.
Conjecture~\ref{conj:ULW} is a global statement
supported by numerical evidence.

%------------------------------
\begin{proposition}[Uniform Spectral Gap]
\label{prop:gap}
There exist $\kappa_0 > 0$ and $c_0 > 0$ such that
for all $c \ge c_0$ and all $l$ with
$|l - \NSh| \le A_0 c^{1/3}$:
\begin{equation}\label{eq:gapbound}
  \lambda_l - \lambda_{l+1} \ge \kappa_0\,(c/2)^{-1/3}.
\end{equation}
\end{proposition}

\begin{remark}[Role and status of Proposition~\ref{prop:gap}]
\label{rem:gap_route}
Proposition~\ref{prop:gap} is the only spectral
condition required by Lemma~\ref{lem:F} and
Corollary~\ref{cor:C}.
It asserts a uniform lower bound on eigenvalue
increments at scale $(c/2)^{-1/3}$, without
specifying the leading-order form.

One route to~\eqref{eq:gapbound} is via
Proposition~\ref{prop:U}(U):
for $c$ sufficiently large, the Airy approximation
implies $\kappa_0 = \tfrac{1}{2}c_{\min}$.
Alternative routes via first-order WKB or
monotone comparison arguments are possible and
would not require the full Airy machinery.
\end{remark}

%------------------------------
\begin{proposition}[ORX Airy Approximation]
\label{prop:U}
There exist $C_1, C_2, c_0 > 0$ (depending only on $A_0$)
such that for all $c \ge c_0$ and all $l$ with
$|l-\NSh| \le A_0 c^{1/3}$:
\begin{itemize}
\item[\textup{(U)}] $\abs{\lambda_l - \FAi(x_l)}
  \le C_1 c^{-2/3}$.
\item[\textup{(U')}] $\abs{(\lambda_l-\lambda_{l+1})
  -(\FAi(x_l)-\FAi(x_{l+1}))}
  \le C_2 c^{-2/3}$.
\end{itemize}
Both bounds are uniform in $l$.
Proposition~\ref{prop:U}\textup{(U)} implies
Proposition~\ref{prop:gap} with
$\kappa_0 = \tfrac{1}{2}c_{\min}$
for $c \ge c_1 := (2C_R/c_{\min})^3$.
\end{proposition}

\begin{remark}[Scope of Proposition~\ref{prop:U}]
\label{rem:propU_scope}
Proposition~\ref{prop:U} is a \emph{local} statement:
it controls individual increments and pointwise values
at the Airy scale $(c/2)^{-1/3}$.
It makes no claim about the global shape of the
eigenvalue profile.
The global erfc structure is addressed separately
in Conjecture~\ref{conj:ULW}.
\end{remark}

\begin{remark}[Why (U') is necessary for Lemma~\ref{lem:A}]
\label{rem:propU_uprime}
Lemma~\ref{lem:A} requires bounding
$E_l - E_{l+1}$ where $E_l := \lambda_l - \FAi(x_l)$.
Condition~(U) gives only $|E_l|, |E_{l+1}| \le C_1 c^{-2/3}$
individually and does not prevent oscillations at scale
$c^{-2/3}$.
Condition~(U') controls the discrete derivative of the
spectral curve and is a genuinely stronger condition.
\end{remark}

\begin{remark}[Route to proof of Proposition~\ref{prop:U}]
\label{rem:propU_route}
ORX \cite{OsipovRokhlinXiao} provides~(U) pointwise.
A proof of both (U) and (U') is expected from:
\begin{enumerate}[(i)]
\item Uniform WKB control with $O(c^{-2/3})$ remainder.
\item Airy matching with uniform remainder; yields~(U).
\item Stability of the spectral map under discrete
  differentiation; yields~(U').
\end{enumerate}
Step~(iii) is the additional content beyond~(U).
\end{remark}

%------------------------------
\begin{conjecture}[Landau--Widom Global Scaling Law]
\label{conj:ULW}
Let $S(c) = (\log c/\pi)^{2/3}$.
There exists a centering parameter $\beta(c) \in \mathbb{R}$
such that for all $l$ in the transition region,
\[
  \lambda_l = \frac{1}{2}
  \operatorname{erfc}\!\left(
  \frac{\ln 2 \cdot (l - \NSh - \beta(c))}{S(c)}
  \right)
  + O(S(c)^{-1}).
\]
Equivalently, the first differences satisfy the
Gaussian boundary layer law
\[
  \lambda_l - \lambda_{l+1}
  \approx \frac{\ln 2}{\sqrt{\pi}\,S(c)}\,e^{-Z_l^2}.
\]
\end{conjecture}

\begin{remark}[Numerical evidence for Conjecture~\ref{conj:ULW}]
\label{rem:ULW_numerics}
For $c \in \{50,100,200\}$, direct eigenvalue computation
via Gauss--Legendre discretization ($n=400$) yields
cross-bandwidth residuals
\[
  \sup_{c,c'}\,
  \sup_{Z \in [-2.5,2.5]}
  |\lambda_c(Z) - \lambda_{c'}(Z)| \approx 2.7 \cdot 10^{-2},
\]
decreasing with $c$.
The centering parameter satisfies
\[
  \beta(50)\approx -0.424,\quad
  \beta(100)\approx -0.380,\quad
  \beta(200)\approx -0.359,
\]
suggesting slow convergence to a limit
$\beta(\infty) \in [-0.30,-0.21]$.
No analytic expression for $\beta(c)$ is known.
\end{remark}

\begin{remark}[Scale separation]
\label{rem:scales}
The three layers operate at genuinely separated scales.
Proposition~\ref{prop:gap} and Lemma~\ref{lem:F} act
at the \emph{combinatorial} scale, requiring only that
gaps are bounded below at order $c^{-1/3}$.
Proposition~\ref{prop:U} resolves the \emph{local Airy}
scale $(c/2)^{-1/3}$.
Conjecture~\ref{conj:ULW} describes the \emph{global
Landau--Widom} scale $S(c) \sim (\log c)^{2/3}$.
The ratio $S(c)/(c/2)^{-1/3} \sim c^{1/3}(\log c)^{2/3}
\to \infty$ confirms genuine scale separation between
the local and global layers.
\end{remark}

% ============================================================
\section{Lemma A: Uniform Airy Increment}
\label{sec:lemmaA}
% ============================================================

\begin{lemma}[Uniform Airy Increment]
\label{lem:A}
Under Proposition~\ref{prop:U}\textup{(U')},
for all $c \ge c_0$ and all $l$ with
$|l-\NSh| \le A_0 c^{1/3}$:
\[
  \lambda_l - \lambda_{l+1}
  = (c/2)^{-1/3}\,\mathrm{Ai}(x_l)^2 + R_l,
\]
where
\[
  |R_l| \le C_R\,c^{-2/3},
  \qquad C_R := L_{A_0}\cdot 2^{2/3} + C_2,
\]
depending only on $A_0$ and $C_2$, uniform in $l$.
Consequently:
\begin{equation}\label{eq:incr_bound}
  |\lambda_l-\lambda_{l+1}|
  \le C_{\max}(c/2)^{-1/3} + C_R\,c^{-2/3}.
\end{equation}
In particular:
\begin{equation}\label{eq:incr_simple}
  |\lambda_l-\lambda_{l+1}|
  \le C_I\,(c/2)^{-1/3},
  \qquad C_I := C_{\max} + C_R\cdot 2^{-1/3}.
\end{equation}
\end{lemma}

\begin{proof}
Write:
\[
  \lambda_l-\lambda_{l+1}
  = [\FAi(x_l)-\FAi(x_{l+1})]
  + [(\lambda_l-\lambda_{l+1})-(\FAi(x_l)-\FAi(x_{l+1}))].
\]
\textbf{Main term.}
$\FAi'=-\mathrm{Ai}^2$; MVT gives
$\FAi(x_l)-\FAi(x_{l+1})
= \mathrm{Ai}(\xi_l)^2(c/2)^{-1/3}$
for some $\xi_l \in [x_l,x_{l+1}]$.
By Lipschitz continuity of $\mathrm{Ai}^2$:
$|\mathrm{Ai}(\xi_l)^2-\mathrm{Ai}(x_l)^2|
\le L_{A_0}(c/2)^{-1/3}$,
contributing $L_{A_0}2^{2/3}c^{-2/3}$ to $R_l$.

\textbf{Increment error.}
By~(U'): $|(\lambda_l-\lambda_{l+1})-(\FAi(x_l)-\FAi(x_{l+1}))|
\le C_2 c^{-2/3}$.

\textbf{Combining:}
$|R_l| \le (L_{A_0}2^{2/3}+C_2)c^{-2/3} = C_R c^{-2/3}$.
Bounds~\eqref{eq:incr_bound} and~\eqref{eq:incr_simple}
follow since $c^{-2/3} \le (c/2)^{-1/3}\cdot 2^{-1/3}$
for $c\ge 1$.
\end{proof}

% ============================================================
\section{Lemma B: Airy Ratio Bound}
\label{sec:lemmaB}
% ============================================================

\begin{lemma}[Airy Ratio Bound]
\label{lem:B}
$\mathrm{Ai}(x)^2/\FAi(x) \le C_{\max}/c_*$
for all $x \in [-A_0,A_0]$.
\end{lemma}

\begin{proof}
Numerator $\le C_{\max}$, denominator $\ge c_*>0$.
\end{proof}

% ============================================================
\section{Lemma F: Dyadic Separation Principle}
\label{sec:lemmaF}
% ============================================================

\begin{lemma}[Dyadic Separation Principle]
\label{lem:F}
Under Proposition~\ref{prop:gap},
for all $c \ge c_0$ and all $i \ne j$
in the transition window:
\begin{equation}\label{eq:specsep}
  |\lambda_i-\lambda_j|
  \ge \kappa_0\,(c/2)^{-1/3}\,|i-j|.
\end{equation}
\end{lemma}

\begin{proof}
Without loss of generality $i < j$.
Since $\{\lambda_l\}$ is strictly decreasing,
by telescoping and Proposition~\ref{prop:gap}:
\[
  |\lambda_i - \lambda_j|
  = \sum_{l=i}^{j-1}(\lambda_l - \lambda_{l+1})
  \ge (j-i)\,\kappa_0\,(c/2)^{-1/3}.\qquad\qedhere
\]
\end{proof}

\begin{remark}[Full independence of Lemma~\ref{lem:F}]
\label{rem:Frobust}
Lemma~\ref{lem:F} depends only on
Proposition~\ref{prop:gap} — not on Airy functions,
erfc approximations, or the specific form of
$\lambda_l - \lambda_{l+1}$.
Any mechanism verifying the gap lower bound~\eqref{eq:gapbound}
suffices: Airy (Proposition~\ref{prop:U}),
WKB-only arguments, or in the future
Conjecture~\ref{conj:ULW} if proved.
\end{remark}

% ============================================================
\section{Corollary C: H2 in Truncated Dyadic Sense}
\label{sec:H2}
% ============================================================

\begin{corollary}[H2 in truncated dyadic sense]
\label{cor:C}
Fix $m_0\ge 1$; let $c_2:=\max(c_0,c_1)$.
Under Propositions~\ref{prop:gap} and~\ref{prop:U}, define
\[
  \delta_m
  := \frac{2C_I}{\kappa_0\cdot 2^{1/3}}\,2^{-m}
  + \frac{2C_{\max}}{c_*}(c/2)^{-1/3}
  + \frac{2C_R}{c_*}c^{-2/3}.
\]
For $c \ge c_2$, $i \le N$, $m \ge m_0$,
$j,j+1 \in \Bm(i)$:
\[
  \frac{|A_{ij}-A_{i,j+1}|}{A_{ij}}
  \le \delta_m = O(2^{-m})+O(c^{-1/3}).
\]
For any $\varepsilon>0$: choose $m_0$ large and
then $c$ large; $\sup_{m\ge m_0}\delta_m < \varepsilon$.
\end{corollary}

\begin{proof}
\[
\frac{|A_{ij}-A_{i,j+1}|}{A_{ij}}
\lesssim
\underbrace{\frac{|\lambda_j-\lambda_{j+1}|}
            {|\lambda_i-\lambda_j|}}_{\text{Term 1}}
+
\underbrace{\frac{|\lambda_j-\lambda_{j+1}|}
            {1-\lambda_j}}_{\text{Term 2}}.
\]

\textbf{Term 1 (Dyadic Separation).}
Numerator $\le C_I(c/2)^{-1/3}$ by
Lemma~\ref{lem:A}\,\eqref{eq:incr_simple}.
Denominator $\ge \kappa_0\,(c/2)^{-1/3}\cdot 2^m$
by Lemma~\ref{lem:F}.
The factor $(c/2)^{-1/3}$ \emph{cancels exactly}:
\[
  \text{Term 1} \le \frac{C_I}{\kappa_0\cdot 2^{-1/3}}\,2^{-m},
  \quad\text{independent of }c.
\]

\textbf{Term 2 (Global error).}
Using~\eqref{eq:incr_bound} and~\eqref{eq:onemlam}:
\[
  \frac{|\lambda_j-\lambda_{j+1}|}{1-\lambda_j}
  \le
  \frac{2C_{\max}}{c_*}(c/2)^{-1/3}
  + \frac{2C_R}{c_*}c^{-2/3} = O(c^{-1/3}).
\]
\end{proof}

\begin{remark}[The operative cancellation]
\label{rem:dyadic}
The cancellation of $(c/2)^{-1/3}$ in Term~1 is
the operative content:
both $|\lambda_j-\lambda_{j+1}|$ and $|\lambda_i-\lambda_j|$
are $O(c^{-1/3})$, but the denominator carries
the combinatorial factor $2^m$ from $|i-j|\ge 2^m$.
Result: purely combinatorial $O(2^{-m})$, independent of~$c$.
The $O(c^{-1/3})$ in $\delta_m$ comes from Term~2 only.
\end{remark}

% ============================================================
\section{Centering Parameter $\beta(c)$}
\label{sec:beta}
% ============================================================

\begin{remark}[Empirical behavior of $\beta(c)$]
\label{rem:beta}
The centering parameter $\beta(c)$ in
Conjecture~\ref{conj:ULW} is determined numerically
by median alignment: $\beta(c)$ is chosen so that
the erfc profile crosses $\lambda = 1/2$ at the
correct index.

Numerical computation via Gauss--Legendre
discretization ($n=400$) yields:
\[
  \beta(50) \approx -0.424,\quad
  \beta(100) \approx -0.380,\quad
  \beta(200) \approx -0.359.
\]
Fitting suggests $\beta(c) \sim A + B/\log(c)$ with
limit $\beta(\infty) \in [-0.30, -0.21]$,
but the model sensitivity is high and no analytic
expression is known.

The observed drift is consistent with a
discretization-induced centering effect, but no
structural identification of $\beta(c)$ in terms
of spectral or arithmetic data is known.
\end{remark}

% ============================================================
\section{Lemma E: Core--WKB Split}
\label{sec:split}
% ============================================================

\begin{lemma}[Core--WKB Split]
\label{lem:E}
Fix $m_0 \ge 1$; assume Propositions~\ref{prop:gap}
and~\ref{prop:U} and B-strong \cite{paperVI}
($A_{ij} \le C_2$).
Then $\sum_{j\ne i}|T_{ij}| \le C(m_0,A,C_2)$
uniformly in $i$, $N$, $c \ge c_2$.
\end{lemma}

\begin{proof}
\textbf{Core ($m<m_0$).}
$|T_{ij}| \le C_2/|i-j|$ (B-strong +
off-diagonal structure, \cite[Lemma~5.2]{paperVII});
sum: $\le 2C_2 m_0$.

\textbf{WKB ($m\ge m_0$).}
Corollary~\ref{cor:C} gives H2.
Apply Theorem~2.1 of \cite{paperVII} to truncated
operator, using H1 and H3: $O(1)$.
\end{proof}

% ============================================================
\section{Conditional Reduction}
\label{sec:closure}
% ============================================================

\begin{corollary}[Conditional Reduction]
\label{cor:closure}
Under Propositions~\ref{prop:gap} and~\ref{prop:U}
and B-strong,
hypotheses (H1)--(H3) of Theorem~2.1 in \cite{paperVII}
are verified, giving
$\sup_{i\le N}\sum_{j\ne i}|P_{ij}|
\le Cc^{-1/2}\log c$
and $\norm{D\mathcal{T}_c^{(N)}} < 1$ for $c\ge c_2$.
\end{corollary}

\begin{remark}[Dependency chain]
\begin{align*}
  &\text{Prop.~\ref{prop:gap}}
    \Rightarrow\text{Lem.~\ref{lem:F}}
    \Rightarrow\text{Cor.~\ref{cor:C} (Term 1)};\\
  &\text{Prop.~\ref{prop:U}(U')}
    \Rightarrow\text{Lem.~\ref{lem:A}}
    \Rightarrow\text{Cor.~\ref{cor:C} (Term 2) + upper bound};\\
  &\text{Prop.~\ref{prop:U}(U)}
    \Rightarrow\text{Prop.~\ref{prop:gap} (one route)};\\
  &\text{Cor.~\ref{cor:C} + H3}
    \Rightarrow\text{Lem.~\ref{lem:E}}
    \Rightarrow B' \Rightarrow A.
\end{align*}
\textbf{Key structural point:}
Lemma~\ref{lem:F} depends only on
Proposition~\ref{prop:gap}.
Proposition~\ref{prop:U} provides the Airy upper bounds
needed for Term~2 in Corollary~\ref{cor:C}
and one route to Proposition~\ref{prop:gap},
but it is not the logical foundation of
the Dyadic Separation Principle.
All thresholds depend only on $A_0, C_1, C_2, c_{\min}$.
\end{remark}

\begin{remark}[Open problems]
\label{rem:open}
\begin{itemize}
\item \textbf{Proposition~\ref{prop:gap}:}
  uniform gap lower bound; plausible from WKB
  or monotone comparison; no proof available.
\item \textbf{Proposition~\ref{prop:U}:}
  the Airy refinement; requires full uniform
  WKB-to-Airy matching with $O(c^{-2/3})$ error control.
\item \textbf{B-strong:} $P_{kl}\le C_2 c^{1/2}$
  in the transition zone; plausible from Widom +
  WKB-to-Airy matching; no proof available.
\item \textbf{Conjecture~\ref{conj:ULW}:}
  a full asymptotic derivation of the erfc scaling law
  and the determination of $\lim_{c\to\infty}\beta(c)$.
\end{itemize}
\end{remark}

% ============================================================
\begin{thebibliography}{9}
\bibitem{DLMF}
{\em NIST Digital Library of Mathematical Functions},
F.~W.~J.~Olver et al.\ (eds.),
Release 1.1.12, \url{https://dlmf.nist.gov/}, 2023.

\bibitem{paperVI}
[Author], \textit{Galerkin Norm Estimates for the Prolate
Phase Operator}, Paper VI in this series.

\bibitem{paperVII}
[Author], \textit{Dyadic Cancellation for PSWF Galerkin
Operators}, Paper VII in this series.

\bibitem{OsipovRokhlinXiao}
A.~Osipov, V.~Rokhlin, H.~Xiao,
\textit{Prolate Spheroidal Wave Functions of Order Zero},
Springer, 2013.
\end{thebibliography}

\end{document}
